% Lösungen Einheit 2 von 3 – Schnittpunkt zweier Geraden (zusammenbau v0.1)
\einheitenkopf{Einheit 2 von 3 · Schnittpunkt zweier Geraden}

\erg{13}{a) drei Gleichungen mit zwei Unbekannten – eine je Koordinate, unbekannt sind nur $r$ und $s$. \quad b) Gleichsetzen: $2 + r = 3$, $2 = s$, $3 + 2r = 3 + s$; aus zwei Gleichungen $r = 1$, $s = 2$; die dritte stimmt als Probe; Schnittpunkt $S(3|2|5)$ \quad c) $h: \vec{x} = (3|0|1) + r \cdot (1|1|b - 1)$; gleichsetzen mit verschiedenen Buchstaben: aus zwei Koordinaten $r = 4$ und $s = 3$; die dritte liefert $b = 2{,}5$ \quad d) Kante $\vec{x} = (0|0|1{,}5) + t \cdot (6|4|z_U - 1{,}5)$, $RT$: $\vec{x} = (2|3|2) + \mu \cdot (2|-2|0)$; aus der ersten und zweiten Koordinate $t = 0{,}5$ und $\mu = 0{,}5$; die dritte Koordinate liefert $z_U = 2{,}5$ – nicht die Höhe der Kante $RT$. \quad e) Die beiden Terme sind die Geraden durch $D$ und $A$ (Richtungsvektor $\overrightarrow{DA}$) und durch $E$ und $B$ (Richtungsvektor $\overrightarrow{EB}$), also die Seitenkanten $DA$ und $EB$. Das Gleichsetzen sucht ihren gemeinsamen Punkt; das Gleichungssystem liefert $r = s = 2$, das Einsetzen in die Kante $DA$ die Spitze $S(-3|4|-6)$ der Pyramide $DEFS$. \quad f) Kante $GH$: $\vec{x} = (0|4|4) + w \cdot (4|0|0)$ mit $0 \le w \le 1$. Gleichsetzen: die zweite Koordinate liefert $u = 1$, die dritte $r^2 = 4$, also $r = 2$ oder $r = -2$; die erste liefert für $r = 2$ den Wert $w = 0{,}25$, für $r = -2$ den Wert $w = 1{,}25$ – außerhalb der Kante, verworfen. Der Schnittpunkt $Q(1|4|4)$ teilt $HG$ im Verhältnis $HQ : QG = 0{,}25 : 0{,}75 = 1 : 3$.}
\erg{14}{a) Lichtgerade $\vec{x} = (0|0|7) + r \cdot (4|-1|-2)$ und Kantengerade $\vec{x} = (6|-3|4) + s \cdot (0|6|0)$ gleichsetzen: $r = 1{,}5$, $s = 0{,}25$, die dritte Gleichung stimmt; Schattenpunkt $(6|-1{,}5|4)$ – er liegt auf der Kante, weil $s$ zwischen 0 und 1 liegt.}
\erg{15}{a) Beide Geraden haben denselben Parameterbuchstaben bekommen, die Parameter sind aber unabhängig. Richtig mit $r$ und $s$: die zweite Koordinate liefert $r = 2$, die dritte $s = 1$; Probe $1 + 2 = 2 + 1$ stimmt; Schnittpunkt $S(3|2|2)$. \quad b) Unbekannt sind nur die beiden Parameter, zwei Gleichungen legen sie fest; die dritte ist eine zusätzliche Bedingung – geht sie nicht auf, gibt es keinen gemeinsamen Punkt.}
